\section{Corner transport and commuting-overlap decompositions}

\begin{theorem}[Transport between spatial compressions]%
    \label{thm:corner_compression_partial_isometry_transport}
    \lean{Matrix.exists_partial_isometry_transporting_compressions}
    \leanok
    Let $P,Q\in\MN D$ be matrices, and suppose that
    $V_P,V_Q:\mathbb C^n\to\mathbb C^D$ satisfy
    \[
        V_P^*V_P=I,\qquad V_PV_P^*=P,\qquad
        V_Q^*V_Q=I,\qquad V_QV_Q^*=Q.
    \]
    Write $\Phi_P(X)=V_PXV_P^*$ and $\Phi_Q(X)=V_QXV_Q^*$ for the
    corresponding linear isomorphisms onto the two corners. For every unitary
    $W\in\MN n$, there is a matrix $U\in\MN D$ such that
    \[
        U=PUQ,\qquad U^*U=Q,\qquad UU^*=P,
    \]
    and, for every $X\in\MN n$,
    \[
        \Phi_P(WXW^*)=U\Phi_Q(X)U^*.
    \]
\end{theorem}

\begin{proof}\leanok
    Take $U=V_PWV_Q^*$. The four isometry identities give the three support
    identities by direct multiplication. Substitution of the two formulas
    for $\Phi_P$ and $\Phi_Q$ gives
    \[
        U\Phi_Q(X)U^*
        =V_PW(V_Q^*V_Q)X(V_Q^*V_Q)W^*V_P^*
        =\Phi_P(WXW^*).
    \]
\end{proof}

\begin{theorem}[Spatial implementation of an isomorphism between corners]%
    \label{thm:corner_star_isomorphism_partial_isometry}
    \lean{Matrix.exists_partial_isometry_implementing_corner_linearEquiv}
    \leanok
    Let $P,Q\in\MN D$ be orthogonal projections. If
    \[
        \Psi:Q\MN DQ\longrightarrow P\MN DP
    \]
    is a bijective complex-linear map preserving multiplication and the
    adjoint, then there is a matrix $U\in\MN D$ such that
    \[
        U=PUQ,\qquad U^*U=Q,\qquad UU^*=P,
    \]
    and
    \[
        \Psi(X)=UXU^*
    \]
    for every $X\in Q\MN DQ$.
\end{theorem}

\begin{proof}\leanok
    \uses{thm:corner_compression_partial_isometry_transport}
    Choose spatial compressions
    \[
        \Phi_Q:\MN{n_Q}\simeq Q\MN DQ,\qquad
        \Phi_P:\MN{n_P}\simeq P\MN DP.
    \]
    The composite $\Phi_P^{-1}\Psi\Phi_Q$ is a linear isomorphism, so equality
    of dimensions gives $n_Q^2=n_P^2$, and hence $n_Q=n_P$. If this common
    rank is zero, both projections vanish and the conclusion holds with
    $U=0$. Otherwise the composite is a star-algebra automorphism of
    $\MN{n_Q}$. The unitary Skolem--Noether theorem writes it as
    $X\mapsto WXW^*$ for a unitary $W$. Applying
    Theorem~\ref{thm:corner_compression_partial_isometry_transport} to $W$
    gives the required matrix $U$.
\end{proof}
\begin{theorem}[Middle-factor coefficients of commuting overlapping operators]%
    \label{thm:commuting_overlapping_middle_coefficients}
    \lean{Matrix.middleSlices_commute_of_overlappingLifts_commute}
    \lean{Matrix.leftOverlappingLift}
    \lean{Matrix.rightOverlappingLift}
    \lean{Matrix.leftMiddleSlice}
    \lean{Matrix.rightMiddleSlice}
    \leanok
    Let $X_{AB}\in\mathbf L(H_A\otimes H_B)$ and
    $Y_{BC}\in\mathbf L(H_B\otimes H_C)$. For fixed outer indices, define
    their middle-factor coefficient matrices by
    \begin{align}
        (X_{aa'})_{bb'}
         & :=(X_{AB})_{(a,b),(a',b')}, \notag \\
        (Y_{cc'})_{bb'}
         & :=(Y_{BC})_{(b,c),(b',c')}.
        \notag
    \end{align}
    If
    \begin{align}
        (X_{AB}\otimes\Id_C)(\Id_A\otimes Y_{BC})
         & =(\Id_A\otimes Y_{BC})(X_{AB}\otimes\Id_C),
        \label{eq:spectral_overlap_commute}
    \end{align}
    then $X_{aa'}Y_{cc'}=Y_{cc'}X_{aa'}$ for every $a,a',c,c'$. This is
    the coefficientwise commutation step in the proof of
    \cite[Lemma~2.1]{Beigi2012Classification}.
\end{theorem}

\begin{proof}
    \leanok
    Evaluate (\ref{eq:spectral_overlap_commute}) between the basis vectors
    $\ket{a,b,c}$ and $\ket{a',b',c'}$. The identity matrices collapse the
    sums over the first and third intermediate indices, leaving
    \begin{align}
        \sum_r(X_{aa'})_{br}(Y_{cc'})_{rb'}
         & =\sum_r(Y_{cc'})_{br}(X_{aa'})_{rb'}.
        \notag
    \end{align}
    This is the $(b,b')$ entry of
    $X_{aa'}Y_{cc'}=Y_{cc'}X_{aa'}$.
\end{proof}


\begin{theorem}[Coefficient form of the commuting-overlap spatial decomposition]%
    \label{thm:commuting_overlapping_coefficient_spatial_decomposition}
    \lean{Matrix.exists_middle_spatial_decomposition_of_overlappingLifts_commute}
    \lean{Matrix.star_rightMiddleSlice}
    \leanok
    Let $X_{AB}\in\mathbf L(H_A\otimes H_B)$ and
    $Y_{BC}\in\mathbf L(H_B\otimes H_C)$, with $Y_{BC}$ Hermitian, and
    suppose $[X_{AB}\otimes\Id_C,\Id_A\otimes Y_{BC}]=0$. There is an
    orthonormal decomposition
    $H_B\cong\bigoplus_{q=0}^{K-1}H_{q,r}\otimes H_{q,l}$ such that all
    middle-factor coefficients $X_{aa'}$ act only on $H_{q,l}$, while all
    middle-factor coefficients $Y_{cc'}$ act only on $H_{q,r}$. More
    precisely, in an orthonormal basis $(e_{q,r,s})$ adapted to this
    decomposition,
    \begin{align}
        X_{aa'}e_{q,r,s}
         & =\sum_{s'}(B^{aa'}_q)_{s's}e_{q,r,s'}, \notag \\
        Y_{cc'}e_{q,r,s}
         & =\sum_{r'}(C^{cc'}_q)_{r'r}e_{q,r',s}.
        \label{eq:spectral_overlap_coeff_actions}
    \end{align}
    The source lemma assumes that both overlapping operators are Hermitian;
    the coefficient conclusion requires Hermiticity only for $Y_{BC}$.
    This is the middle-space coefficient form of
    \cite[Lemma~2.1]{Beigi2012Classification}.
\end{theorem}

\begin{proof}
    \leanok
    \uses{thm:commuting_overlapping_middle_coefficients}
    Let $S$ be the commutant of the family $(Y_{cc'})_{c,c'}$ together with
    its adjoints. Hermiticity gives $Y_{cc'}^*=Y_{c'c}$.
    Theorem~\ref{thm:commuting_overlapping_middle_coefficients} therefore
    places every $X_{aa'}$ in $S$. Conversely, every $Y_{cc'}$ commutes
    with all members of $S$. Apply the complementary spatial-actions
    Theorem~\ref{thm:starSubalgebra_complementary_spatial_actions}
    to $S$. Its two action formulas, specialized to $X_{aa'}$ and $Y_{cc'}$,
    give (\ref{eq:spectral_overlap_coeff_actions}).
\end{proof}

\begin{theorem}[Explicit unitary block actions for commuting overlaps]%
    \label{thm:commuting_overlapping_unitary_block_actions}
    \lean{Matrix.exists_unitary_blockActions_of_overlappingLifts_commute}
    \lean{Matrix.leftSpatialBlockEquiv}
    \lean{Matrix.rightSpatialBlockEquiv}
    \leanok
    Let $X_{AB}\in\mathbf L(H_A\otimes H_B)$ and
    $Y_{BC}\in\mathbf L(H_B\otimes H_C)$ be Hermitian, and suppose
    $[X_{AB}\otimes\Id_C,\Id_A\otimes Y_{BC}]=0$. There are positive
    integers $d_q,m_q$, a unitary identification
    $H_B\cong\bigoplus_{q=0}^{K-1}H_{q,l}\otimes H_{q,r}$, and Hermitian
    operators $R_q\in\mathbf L(H_A\otimes H_{q,l})$ and
    $S_q\in\mathbf L(H_{q,r}\otimes H_C)$ such that
    \begin{align}
        (\Id_A\otimes U^\dagger)X_{AB}(\Id_A\otimes U)
         & =\bigoplus_{q=0}^{K-1}
        (R_q\otimes\Id_{H_{q,r}}), \notag \\
        (U^\dagger\otimes\Id_C)Y_{BC}(U\otimes\Id_C)
         & =\bigoplus_{q=0}^{K-1}
        (\Id_{H_{q,l}}\otimes S_q).
        \label{eq:spectral_overlap_block_actions}
    \end{align}
    This is the explicit coordinate form of the two block-action identities
    in~\cite[Lemma~2.1]{Beigi2012Classification}.
\end{theorem}

\begin{proof}
    \leanok
    \uses{thm:commuting_overlapping_coefficient_spatial_decomposition}
    Use the orthonormal decomposition from
    Theorem~\ref{thm:commuting_overlapping_coefficient_spatial_decomposition}.
    Let $U$ be the unitary whose columns are the adapted basis vectors.
    Assemble the coefficient matrices $(B_q^{aa'})_{a,a'}$ into a matrix
    $R_q$ on $H_A\otimes H_{q,l}$, and similarly assemble
    $(C_q^{cc'})_{c,c'}$ into a matrix $S_q$ on
    $H_{q,r}\otimes H_C$. Taking matrix entries in the adapted basis turns
    (\ref{eq:spectral_overlap_coeff_actions}) into
    (\ref{eq:spectral_overlap_block_actions}). Unitary conjugation preserves
    Hermiticity. Restricting the two conjugated Hermitian operators to a
    summand, and fixing one basis vector in the nonempty complementary
    factor, shows that every $R_q$ and every $S_q$ is Hermitian.
\end{proof}

\begin{theorem}[Block-sum equalities for commuting overlapping operators]%
    \label{thm:commuting_overlapping_block_sum_equalities}
    \lean{Matrix.exists_unitary_blockSum_equalities_of_overlappingLifts_commute}
    \leanok
    \uses{thm:commuting_overlapping_unitary_block_actions}
    Let $X_{AB}\in\mathbf L(H_A\otimes H_B)$ and
    $Y_{BC}\in\mathbf L(H_B\otimes H_C)$ be Hermitian, and suppose
    $[X_{AB}\otimes\Id_C,\Id_A\otimes Y_{BC}]=0$. There are positive
    integers $d_q,m_q$, an identification
    $H_B\cong\bigoplus_{q=0}^{K-1}H_{q,l}\otimes H_{q,r}$, a unitary
    $U\in\mathbf L(H_B)$, and Hermitian operators
    $R_q\in\mathbf L(H_A\otimes H_{q,l})$ and
    $S_q\in\mathbf L(H_{q,r}\otimes H_C)$ such that
    \begin{align}
        X_{AB}
         & =(\Id_A\otimes U)
        \left(\bigoplus_{q=0}^{K-1}
        R_q\otimes\Id_{H_{q,r}}\right)
        (\Id_A\otimes U^\dagger), \notag \\
        Y_{BC}
         & =(U\otimes\Id_C)
        \left(\bigoplus_{q=0}^{K-1}
        \Id_{H_{q,l}}\otimes S_q\right)
        (U^\dagger\otimes\Id_C).
        \label{eq:spectral_overlap_block_sums}
    \end{align}
    In these formulas the two direct sums are transported to the original
    product coordinates through the stated identification. These are the
    two equalities in~\cite[Lemma~2.1]{Beigi2012Classification}. The
    operators $R_q$ and $S_q$ are Hermitian; they are not asserted to be
    unitary.
\end{theorem}

\begin{proof}
    \leanok
    Let $V_A=\Id_A\otimes U$ and $V_C=U\otimes\Id_C$.
    Theorem~\ref{thm:commuting_overlapping_unitary_block_actions} gives
    \begin{align}
        V_A^\dagger X_{AB}V_A
         & =\bigoplus_qR_q\otimes\Id_{H_{q,r}}, \notag \\
        V_C^\dagger Y_{BC}V_C
         & =\bigoplus_q\Id_{H_{q,l}}\otimes S_q.
        \notag
    \end{align}
    Since $V_AV_A^\dagger=\Id$ and $V_CV_C^\dagger=\Id$, conjugating these
    identities by $V_A$ and $V_C$, respectively, gives
    (\ref{eq:spectral_overlap_block_sums}).
\end{proof}
